\section{Study Population}

The study population comprises adults with histologically confirmed, resectable or
borderline-resectable, KRAS G12-mutated pancreatic ductal adenocarcinoma (PDAC)
who are eligible for perioperative daraxonrasib and whose anatomy is compatible
with the eight-arm robotic system. Target enrollment is approximately 36
participants screened to yield 18 treated, consistent with the sample size in
\S9. Figure~6 presents the CONSORT-style participant flow from screening through
the 3+3 dose-escalation cohorts to the analysis populations.

\begin{mermaidfig}
\node[mmin]  (as) at (0,0)      {Assessed for\\eligibility\\(n approx.\ 36\\screened)};
\node[mmdark] (ex) at (6.9,0)    {Excluded: not meeting \S5.1 / meeting \S5.2; declined;\\unresectable on staging};
\node[mmstep](en) at (0,-3.4)   {Enrolled and\\dose-assigned\\(n = 18 treated,\\3+3 design)};
\node[mmstep](d1) at (4.6,-3.4) {DL1 cohort\\n = 3 (expand to\\6 on 1 DLT)};
\node[mmstep](d2) at (9.2,-3.4) {DL2 cohort\\n = 3-6};
\node[mmgoal](tr) at (0,-6.8)   {Received robotic\\Whipple +\\perioperative\\daraxonrasib};
\node[mmstep](d3) at (4.6,-5.4) {DL3 cohort\\(target dose)\\n = 3-6};
\node[mmdark] (di) at (6.9,-8.5) {Discontinued\\intervention\\(conversion,\\withdrawal, \S7)};
\node[mmgoal](an) at (0,-10.0)  {Analyzed: safety\\(all dosed),\\DLT-evaluable,\\per-protocol, mITT};
\draw[mmedge] (as) -- node[above,font=\scriptsize]{exclude} (ex);
\draw[mmedge] (as) -- (en);
\draw[mmedge] (en) -- (d1);
\draw[mmedge] (d1) -- (d2);
\draw[mmedge] (en) -- (tr);
\draw[mmedge] (tr) -- (d3);
\draw[mmedged] (tr) -| node[below,font=\scriptsize,pos=0.25, xshift=0cm, yshift=-0.1cm]{protocol-defined} (di);
\draw[mmedge] (tr) -- (an);
\end{mermaidfig}
\vspace{-0.7cm}
\figcaption{Figure 6. CONSORT-style participant flow from screening
(approximately 36 assessed) through the 3+3 dose-escalation cohorts (DL1 through
DL3, up to n = 18 treated) to the analysis populations, with the\\exclusion and
protocol-defined discontinuation branches that feed the per-protocol and safety
sets.}

\clearpage

\subsection{Inclusion Criteria}

To be eligible to participate in this study, an individual must meet all of the
following criteria.

\begin{enumerate}[leftmargin=1.6em,itemsep=2pt,topsep=2pt]
\item Signed and dated informed consent obtained before any study-specific
procedure, including the participant's explicit election or declination of the
Physical AI opt-out, which governs use of the LLM-directed advisory layer (\S8).
\item Age $\geq 18$ years at the time of consent.
\item Histologically or cytologically confirmed PDAC that is resectable or
borderline-resectable on multidisciplinary radiologic review.
\item A documented KRAS G12 mutation (for example G12D, G12V, or G12R) on a
validated tumor genotyping assay.
\item Eastern Cooperative Oncology Group (ECOG) performance status of 0 or 1.
\item Adequate organ and marrow function on screening laboratories, including
hematologic, hepatic, and renal indices sufficient to undergo major
pancreaticobiliary surgery and receive daraxonrasib.
\item A measured baseline CA 19-9 value (no threshold imposed; value recorded for serial assessment).
\item Eligible for perioperative daraxonrasib under broad pan-KRAS criteria of RASolute 302 \cite{rasolute302}.
\item Tumor and vascular anatomy compatible with the eight-arm robotic system on
the preoperative imaging review, including port placement and working volume
adequate for operative tasks.
\end{enumerate}

The synthetic reference exemplar PAT-PDAC-0001 illustrates a representative
eligible participant: a 68-year-old with a 3.4~cm head-of-pancreas tumor abutting
the superior mesenteric vein at $75^\circ$, KRAS G12D, microsatellite stable,
CA 19-9 of 412~U/mL at diagnosis, ECOG performance status 1, after a completed
neoadjuvant modified FOLFIRINOX course of four cycles \cite{Conroy2018}. This
exemplar is used for design illustration and does not represent an enrolled
individual.

\subsection{Exclusion Criteria}

An individual who meets any one of the following criteria is excluded from
participation.

\begin{enumerate}[leftmargin=1.6em,itemsep=2pt,topsep=2pt]
\item Distant metastatic disease on staging.
\item Vascular involvement that precludes a margin-negative (R0) resection, or
anatomy otherwise incompatible with the eight-arm robotic approach.
\item Prior pancreatic resection or prior surgery that precludes the planned
reconstruction.
\item Uncontrolled intercurrent illness or comorbidity (for example uncontrolled
cardiac, pulmonary, hepatic, renal, or active infectious disease) that would make
major surgery or daraxonrasib administration unsafe.
\item Pregnancy or lactation, or unwillingness to use adequate contraception
during study if applicable.
\item Known contraindication or hypersensitivity to daraxonrasib or to a required
excipient, or a concomitant medication interaction that cannot be safely managed.
\item Inability to provide informed consent or to comply with the Schedule of
Activities.
\end{enumerate}

Participant selection is equitable. Eligibility is determined solely by the
scientific and safety criteria above, and no eligible individual is excluded on
the basis of sex, gender, race, ethnicity, socioeconomic status, or insurance
status. Limited English proficiency is not an exclusion criterion; qualified
interpreters and translated consent materials are provided so that individuals
with limited English proficiency can be fully informed and can participate on the
same basis as English-speaking individuals.

\subsection{Lifestyle Considerations}

Perioperative lifestyle restrictions apply and are reviewed at consent and at the
baseline visit. Participants follow a standardized perioperative nutrition plan,
including the institutional preoperative fasting protocol and a staged
postoperative diet advancement keyed to recovery of gastrointestinal function and
to fistula status. Alcohol is to be avoided during the perioperative period and
during any interval of active daraxonrasib dosing, both to limit hepatotoxic
burden and to avoid confounding the safety assessment. Tobacco cessation is
strongly advised and is supported with referral, given its effect on wound healing
and pulmonary risk. No restriction on routine physical activity is imposed beyond
standard postoperative surgical guidance. There are no study-specific dietary
caffeine or grapefruit restrictions beyond those required by daraxonrasib labeling
and the concomitant-medication rules in \S6.

\subsection{Screen Failures}

A screen failure is an individual who consents to participate but is not
subsequently enrolled because one or more eligibility criteria are not met. For
each screen failure, a minimal CONSORT screen-failure data set is recorded:
demography, the specific criterion or criteria that were not met (the
screen-failure detail), the eligibility determination, and any serious adverse
event that occurred between consent and the determination. No other study
procedures are performed after a screen-failure determination. An individual who
fails screening may be rescreened if the disqualifying condition was transient and
is expected to resolve (for example a correctable laboratory abnormality);
a rescreened individual retains the same participant number, and the
rescreening is documented so that the screening-to-enrollment accounting in
Figure~6 remains complete.

\subsection{Strategies for Recruitment and Retention}

Recruitment is conducted at a single academic medical center through the
institutional multidisciplinary pancreatic-cancer program. Eligible individuals
are identified at tumor-board review and at surgical and medical-oncology clinics,
and the study is described to them by a member of the research team who is not the
treating surgeon, to reduce the potential for undue influence. The planned accrual
rate is approximately one to two participants per month, consistent with the
sentinel staggered enrollment of the 3+3 design and the Phase 0 readiness gate,
giving an enrollment period sufficient to treat 18 participants from approximately
36 screened.

Retention through the 24-month overall-survival window is supported by structured
follow-up at day 30, day 90, and every 12 weeks thereafter, with appointment
reminders, flexible scheduling, coordination with referring providers for
participants who travel, and reimbursement of reasonable travel and parking costs.
Because the population is vulnerable, additional safeguards consistent with Office
for Human Research Protections (OHRP) guidance are applied: enrollment uses an
independent consent monitor when indicated, the consent process explicitly
addresses therapeutic misconception about the LLM-directed robotic system and the
voluntariness of the Physical AI opt-out, and participants are reminded at each
contact that they may withdraw at any time without affecting their clinical care.
Compensation is limited to reimbursement of documented study-related travel,
parking, and meal expenses, provided to the participant by the site after each
completed study visit; no payment is offered for undergoing the surgical procedure
or the investigational drug, so that compensation does not function as an undue
inducement.
